\chapter{Implementation Details}
\label{ch:implementation}

Chapter~\ref{ch:architecture} defined the architectural design objective that governs Antares AI: keeping personally identifiable information out of the LLM context boundary in cleartext. This chapter shows how that objective is pursued in running code. Section~\ref{sec:impl_tech_stack} describes the technology stack and the reasoning behind splitting the system across two language ecosystems. Section~\ref{sec:impl_anonymization_pipeline} covers the two implementation phases of the reversible anonymisation pipeline, from an early Presidio-based prototype to the GLiNER-based system now in use. Section~\ref{sec:impl_mcp_tools} details how business operations are exposed to Claude as MCP tools. Section~\ref{sec:impl_jobtools} describes the orchestration logic behind the three use cases. Section~\ref{sec:impl_backend_integration} closes the chapter with the abstraction layer that separates business logic from the Miorelli backend it currently mocks.

% ==========================================
\section{Technology Stack and Development Environment}
\label{sec:impl_tech_stack}

Antares AI is split across two language ecosystems, with the MCP server, the MCP client, and the Anonymisation Service built on .NET~10 with ASP.NET Core~\cite{dotnet10}, while the GLiNER NER microservice is implemented in Python with FastAPI~\cite{fastapi2018} and served by Uvicorn. This division reflects the specific requirements of GLiNER, which relies on HuggingFace \texttt{transformers} and Python-native inference stacks, tokenisers, and model-loading utilities for which no stable equivalent exists in .NET. Conversely, the remainder of the system leverages C\#, strong typing, and reflection facilities, enabling the MCP server to generate tool schemas from C\# attributes while utilising first-class .NET SDKs for Anthropic and Model Context Protocol integrations. Splitting microservices along runtime capability boundaries rather than enforcing a single language across all services represents standard architectural practice~\cite{newman2021building}, aligning here precisely with the machine learning inference component.

Redis~\cite{redis2009} serves as the sole piece of infrastructure shared indirectly across the ecosystem split, containing no application code and remaining reachable exclusively from the Anonymisation Service, which uses it to store token-to-entity mappings for each context. Because Redis operates as a network-accessible key-value store rather than a shared library or database schema, both ecosystems remain completely decoupled at the code level, meaning the GLiNER service is unaware of Redis and the C\# services interact with it solely through the \texttt{IAnonymizationContextStore} contract detailed in Section~\ref{subsec:impl_token_reinjection}.

Dependency management adheres to ecosystem-specific conventions rather than a unified project standard. The three .NET projects declare packages in \texttt{.csproj} files where \texttt{dotnet restore} executes automatically on build, eliminating manual installation steps. The GLiNER Service declares dependencies in \texttt{pyproject.toml} and utilises \texttt{uv}~\cite{astral2024uv} for resolution and environment management, executing \texttt{uv sync} automatically via \texttt{uv run} on first launch. Both mechanisms ensure developers can clone and execute services without manual setup. The sole point of code-level integration between the two ecosystems occurs via NSwag~\cite{nswag2015}, which reads the OpenAPI specification of the Anonymisation Service to generate a typed C\# client consumed by both the MCP server and the MCP client.

% ==========================================
\section{The Reversible Anonymisation Pipeline}
\label{sec:impl_anonymization_pipeline}

The anonymisation pipeline evolved through two implementation phases. The initial approach utilised Microsoft Presidio, designed prior to concrete use cases, but failed when confronted with real requests. The subsequent architecture centered on GLiNER, replacing Presidio once actual Italian business prompts revealed operational shortcomings. Both phases share a downstream requirement wherein token-to-value mappings must reside in a secure location accessible to the server but inaccessible to the client, as described in Section~\ref{subsec:impl_token_reinjection}.

\subsection{Early Approach: Microsoft Presidio and its Limitations}
\label{subsec:impl_presidio}

Presidio~\cite{microsoft2018presidio} served as the initial detection engine selected during architecture design due to its leading open-source status, predictable recognisers, and local execution without external inference dependencies, as outlined in Section~\ref{subsec:background_presidio_gliner}. Because this selection preceded the acquisition of the three concrete use cases, it relied on general merit rather than benchmarked performance against actual inputs.

Practical limitations emerged rapidly upon evaluating real prompts. Presidio relies on pattern- and rule-based recognisers tuned for English-language, general-purpose PII, failing to match domain vocabulary such as \texttt{cantiere}, \texttt{PULITO\_DA}, and \texttt{Capocantiere}, or the informal spoken Italian register with embedded names characteristic of incoming prompts. Entities identifiable immediately by humans, such as personal names within holiday requests, were silently omitted rather than flagged with low confidence. While Presidio recognisers assign fixed per-pattern confidence scores defined by authors rather than inferred from inputs, the absence of matching patterns results in complete silence rather than partial matches or thresholdable scores. Extending coverage manually for every entity type and variant in Italian operational text proved incompatible with input variability, reflecting the structural limitations discussed in Section~\ref{subsec:background_presidio_gliner} as concrete recall failures on project test prompts.

\subsection{The Pivot to GLiNER for Multilingual Italian NER}
\label{subsec:impl_gliner_pivot}

The anonymisation module was restructured around GLiNER~\cite{zaratiana2023gliner}, a zero-shot NER model accepting arbitrary label sets at inference time instead of fixed pre-trained schemas. The deployed checkpoint, \texttt{urchade/gliner\_multi-v2.1}~\cite{zaratiana2024glinermultiv21}, supports multilingual text natively without fine-tuning for Italian, directly resolving previous recognition limitations through configurable entity types rather than hardcoded recognisers.

The GLiNER Service exposes a single \texttt{/analyze} endpoint accepting text strings, label lists, confidence thresholds, and language hints, returning entity spans with positions and scores. The default label set encompasses eight categories, specifically person, organisation, location, address, email, phone number, date, and time. Configured with a threshold of 0.45 established during initial setup without systematic sweeps, the operating point leans permissive by design, treating missed entities as privacy failures while minor, correctable token annoyances represent acceptable trade-offs in reconstructed text, favoring recall over precision. Determining optimal thresholds represents a studied problem in entity extraction, where deployments calibrate operating points against specific miss versus false alarm costs~\cite{minkov2006tradeoff}. While systematic validation sweeps define exact precision-recall curves, Chapter~\ref{ch:usecases} reports empirical results from this uncalibrated choice without claiming optimality.

Raw GLiNER predictions undergo a three-stage filtering pipeline before reaching callers. The normalisation stage discards predictions with out-of-bounds offsets or empty labels, reconstructing text strings when omitted from raw output. The quality-filtering stage removes entities falling below confidence thresholds, spanning fewer than three characters, or matching Italian stopwords from a project-specific blocklist containing operational imperatives like \emph{lista}, \emph{risolvi}, and \emph{internamente} frequently misidentified as entities. The deduplication stage resolves overlapping spans, such as \emph{Giulia} and \emph{Giulia Conti}, by sorting candidates by score and length, greedily retaining the best non-overlapping set and restoring position order, mirroring greedy suppression rules used for object detection bounding boxes~\cite{neubeck2006nms}. These sequential filters enhance precision following initial recall, achieving low false-negative rates on real entities without generating excessive false positives.

\subsection{Deterministic Token Mapping and Redis Context Storage}
\label{subsec:impl_token_reinjection}

Following entity detection by GLiNER or Presidio, the Anonymisation Service substitutes entities with tokens formatted as \texttt{<LABEL\_NNNN>}, combining uppercase labels with zero-padded four-digit counters scoped per label within the current context. A prompt containing two distinct individuals generates \texttt{<PERSON\_0001>} and \texttt{<PERSON\_0002>}, preserving distinctions for downstream reasoning including LLM evaluation.

Substitution proceeds from right to left by sorting entities by start offset and replacing them backward, preventing length discrepancies from shifting character offsets of subsequent entities and corrupting detection coordinates. The resulting token-to-value mapping is persisted in Redis~\cite{redis2009} under a newly generated context identifier through the \texttt{IAnonymizationContextStore} contract supporting four operations, namely \texttt{Store} for new mappings, \texttt{Get} for non-destructive reads, \texttt{Append} for multi-turn session merges, and \texttt{Pop} for combined read and delete execution. Although \texttt{Pop} was designed for server-side deanonymisation, Section~\ref{subsec:arch_anonymization_service} notes that \texttt{/deanonymize-known} invokes \texttt{Get}, leaving active implementations without code paths removing mappings from Redis. Combined with unconfigured time-to-live properties, context persistence is bounded solely by process lifecycles. This reflects a deliberate proof-of-concept trade-off for local operations lacking persistent volumes, validating end-to-end orchestration rather than production-grade retention policies.

% ==========================================
\section{Exposing Business Operations via MCP}
\label{sec:impl_mcp_tools}

Section~\ref{subsec:background_mcp_standardization} detailed abstract MCP primitives, whereas this section examines how \texttt{Antares\_AI\_McpServer} converts C\# classes into executable tools for Claude, covering schema generation in Section~\ref{subsec:impl_csharp_reflection}, routing entry points in Section~\ref{subsec:impl_ask_antares}, and privacy-enforcing system prompts in Section~\ref{subsec:impl_prompting}.

\subsection{C\# Attributes and Reflection for Tool Schema Generation}
\label{subsec:impl_csharp_reflection}

Classes exposing MCP tools carry \texttt{[McpServerToolType]}, and callable methods carry \texttt{[McpServerTool(Name = "...")]}. The \texttt{ModelContextProtocol.AspNetCore} package inspects these attributes via startup reflection, constructing JSON Schema~\cite{jsonschema2020} tool definitions~\cite{anthropic2024mcp} derived directly from method parameters and names. Developers add business operations by implementing standard C\# methods and attributes, automatically generating tool schemas, HTTP/SSE exposure, and Claude function-integration bindings without manual schema synchronisation.

Only two classes, namely \texttt{AntaresTools} and \texttt{JobTools}, register as MCP tool types via \texttt{.WithTools<T>()} in \texttt{program.cs}, while remaining tool classes operate as scoped services injectable into \texttt{JobTools} but invisible to Claude:
\begin{itemize}
\item \texttt{DepartmentsTools}
\item \texttt{IdentityTools}
\item \texttt{AssetTools}
\item \texttt{PersonTools}
\item \texttt{AbsenceTools}
\item \texttt{LookupTools}
\end{itemize}
This narrows the reachable model surface in accordance with foundational least privilege principles formulated by Saltzer and Schroeder~\cite{saltzer1975protection}, ensuring compromised components cause restricted damage by limiting model agency exclusively to business-level operations while insulating data access layers. Claude performs intent classification and parameter extraction without requiring direct access to data management classes like \texttt{PersonTools} or \texttt{AssetTools}.

\subsection{The \texttt{ask\_antares} Routing Tool}
\label{subsec:impl_ask_antares}

\texttt{AntaresTools.AskAntares} serves as the sole tool invoked by the MCP client, accepting \texttt{user\_input} containing anonymised tokens and optional \texttt{context\_id} parameters for server-side deanonymisation propagation. Internally, it constructs an Italian system prompt designating Claude as a semantic router constrained to select exactly one of five available sub-tools, forward \texttt{context\_id} values unchanged, avoid restating or revealing PII, and return hardcoded confirmation strings instead of freeform summaries. It then invokes \texttt{IChatClient.GetResponseAsync()}, binding C\# calls to Claude's schema-driven tool-use contract~\cite{anthropic2024tooluse} with function invocation enabled, returning raw response text. Executing at a sampling temperature of 1.0, the model output distribution remains unmodified, preserving low-probability tails later addressed in decoding strategies like nucleus sampling~\cite{holtzman2020degeneration}. While this tail introduces potential run-to-run variance evaluated in validation phases, repeated invocations lack deterministic guarantees, constituting a methodological limitation of the study.

The five sub-tools Claude can select from are:
\begin{itemize}
\item \texttt{GetDepartments}
\item \texttt{GetDepartement}
\item \texttt{GetNearestWorkersToCantiere}
\item \texttt{GetColleaguesOfPerson}
\item \texttt{ReportHolidayRequest}
\end{itemize}
The initial two represent early prototyping artifacts, while the latter three address the core use cases examined in Chapter~\ref{ch:usecases}. Claude operations terminate upon selecting sub-tool names and extracting parameters, delegating downstream processing, deanonymisation, GUID resolution, ranking, and report generation to \texttt{JobTools} code invisible to the model beyond final confirmation strings.

Transient Claude API failures trigger retries up to four times using exponential backoff intervals of 2, 4, and 8 seconds~\cite{brooker2015backoff}, returning hardcoded Italian error strings only after total failure to bound blocking behaviour without premature dropping.

\subsection{Strict Prompting for Privacy Enforcement}
\label{subsec:impl_prompting}

System prompts within \texttt{AskAntares} enforce privacy via instructions where structural designs are inapplicable due to boundary conditions. While upstream components including anonymisation pipelines, token substitution, and Redis stores isolate PII per privacy-by-design principles from Section~\ref{sec:arch_privacy_by_design}, Claude operates externally via instruction rather than structural verification. System prompts incorporate three constraints: retaining anonymised tokens without paraphrasing or translation, prohibiting unprovided names, phone numbers, or addresses, and forwarding \texttt{context\_id} attributes to selected sub-tools for server-side deanonymisation.

These instructions function as mitigations rather than formal proofs, remaining susceptible to prompt disregard and indirect injection vulnerabilities outlined in Section~\ref{subsec:background_mcp_security}~\cite{greshake2023not}. Because \texttt{AskAntares} returns model responses unmodified, fabricated PII remains uncaught at boundaries, relying instead on upstream efficacy ensuring models lack sensitive data to misuse. \texttt{JobTools} resolves parameters and queries production data post-selection, returning hardcoded success strings while writing sensitive reports to server-local logs inaccessible to Claude, ensuring successful detection prevents data exposure regardless of instruction adherence.

% ==========================================
\section{Use Case Orchestration in \texttt{JobTools}}
\label{sec:impl_jobtools}

The three sub-tools available for routing by Claude, described in Section~\ref{subsec:impl_ask_antares}, reside within \texttt{JobTools} and follow a consistent five-step structure. First, any natural-language parameter extracted by Claude, such as person or cantiere names, is deanonymized via \texttt{/deanonymize-known}. If \texttt{context\_id} is null or the call fails, parameters pass through unchanged rather than blocking requests, ensuring graceful degradation when upstream anonymization is omitted. Second, the operation resolves required relation-type and category GUIDs, as the Miorelli API utilises GUID addressing exclusively. Third, lower-level Miorelli tools are queried to assemble relevant datasets scoped to cantieri authorised for the caller. Fourth, specific business logic is applied, involving ranking, authorisation checks, or absence cross-referencing. The sequence of steps three and four varies across use cases, as Use Case C, detailed in Section~\ref{subsec:impl_ucc}, executes authorisation checks prior to dataset assembly to avoid querying absence data for unauthorised requests. Fifth, comprehensive reports containing deanonymized data are written to server-local log files, returning identical hardcoded Italian confirmation strings regardless of report contents.

Steps one and five embody the architectural privacy objectives from Chapter~\ref{ch:architecture} in practice. Step one represents the sole reentry point for real data post-anonymisation, while step five represents the exclusive exit point directed towards local disk, ensuring data never flows back towards Claude or clients.

\subsection{Use Case A: Nearest Workers to a Cantiere}
\label{subsec:impl_uca}

\texttt{GetNearestWorkersToCantiere} accepts cantiere names and optional search radii defaulting to 15\,km. Following name deanonymization, it resolves caller identities, \texttt{PULITO\_DA} relation types, and \texttt{SedePrincipale} item categories, invoking \texttt{GetAuthorizedCantieri()} to match deanonymized names against authorised sites. Upon locating target cantieri, \texttt{GetPeopleByRelationOnAsset} retrieves linked workers via \texttt{PULITO\_DA}, fetching individual worker details and residential coordinates.

Ranking workers by distance utilises the Haversine formula~\cite{sinnott1984haversine} implemented in \texttt{CalculateDistanceAndRank} to compute great-circle distances on a sphere from latitude and longitude coordinates. For short-range travel across local workers, spherical formulas account for Earth curvature ignored by flat-plane approximations. Workers lacking coordinates are ranked last, while \texttt{FilterByRadiusKm} excludes those exceeding requested radii. Results are partitioned into \emph{Attivi} and \emph{NonAttivi} before logging full reports containing names, phone numbers, and computed distances. Report generation handles three edge cases generating distinct log entries: unrecognised deanonymized names, unrecognised services or cantieri, and authorised cantieri falling outside caller perimeters. Responses delivered to Claude and clients remain identical across all branches.

\subsection{Use Case B: Colleagues of a Person}
\label{subsec:impl_ucb}

\texttt{GetColleaguesOfPerson} addresses two queries for a deanonymized person simultaneously: identifying direct colleagues and discovering nearby non-colleagues. Following person resolution, \texttt{GetColleaguesOfPerson} queries Miorelli services to return individuals sharing at least one \texttt{PULITO\_DA} cantiere under relational queries. The method then scans authorised cantieri to collect workers lacking shared sites with the target person, ranking non-colleagues by residential distance using \texttt{CalculateDistanceAndRank} before retaining the five closest matches. Logs record both direct colleagues and nearest non-colleagues prior to returning confirmation strings.

\subsection{Use Case C: Holiday Request Analysis}
\label{subsec:impl_ucc}

\texttt{ReportHolidayRequest} incorporates pre-execution authorisation checks to determine eligibility. Following target person deanonymization and relation-type resolution for \texttt{RESPONSABILE\_DI\_RIFERIMENTO} and \texttt{AUTORIZZATORE}, methods verify caller eligibility across three conditions: managing the person within \texttt{PULITO\_DA} perimeters, serving as their Responsabile, or serving as their Autorizzatore. Authorisation executes deterministically in C\# prior to absence queries, terminating Claude involvement once intent selection and parameter extraction conclude.

Upon authorisation, worker pools spanning authorised cantieri are constructed to identify target site assignments. For each cantiere, \texttt{GetAbsencesForPerson} queries site workers to identify absent colleagues, filtering nearby non-colleagues within 15\,km to flag absences. Reports provide per-cantiere breakdowns since substitutes require spatial proximity and operational availability on specific dates. Consistent with other use cases, callers and Claude receive only confirmation strings while detailed breakdowns reside solely in server-side logs.

% ==========================================
\section{Backend Integration and Data Abstraction}
\label{sec:impl_backend_integration}

Use cases detailed in Section~\ref{sec:impl_jobtools} require operational workforce data covering persons, cantieri, relations, and absences. Tool classes retrieving this data bypass direct backend communication, interacting exclusively with \texttt{IMiorelliService}, which operated via mock implementations throughout the internship.

\subsection{The \texttt{IMiorelliService} Interface}
\label{subsec:impl_miorelli_interface}

\texttt{IMiorelliService} serves as the singular interface utilised by data-access tool classes including \texttt{PersonTools}, \texttt{AssetTools}, \texttt{AbsenceTools}, \texttt{LookupTools}, and \texttt{IdentityTools}. These classes avoid HTTP request construction and API schema dependencies, invoking strongly typed interface methods instead. Depending on abstractions rather than concrete backends implements dependency inversion~\cite{martin1996dip}, isolating high-level business logic in \texttt{JobTools} from low-level retrieval details and allowing interchangeable backend implementations without modifying calling code.

Certain tool methods perform internal computations, such as \texttt{CalculateDistanceAndRank} on \texttt{PersonTools}, which executes pure calculations over data returned by \texttt{IMiorelliService}, as the interface supplies base entities rather than rankings.

\subsection{The Stub/Mock Pattern}
\label{subsec:impl_miorelli_stub}

\texttt{MiorelliServiceStub} served as the sole implementation of \texttt{IMiorelliService} during the internship, returning in-memory data from fictional workers and cantieri like Giulia Conti and Bergamo rather than querying production databases.

This scope decision avoided dividing team focus between anonymisation engineering and live sensitive data integration. Prioritising correct architectures over immediate production completion ensured data safety.

The stub also simulates failures by throwing \texttt{MiorelliClientNotAvailableException} for missing records, prompting \texttt{IdentityTools.ResolveCallerIdentity} to catch exceptions and supply default fallback identities, specifically Paolo Ricci, Service Manager, implementing Fowler's Special Case pattern~\cite{fowler2002specialcase} to provide well-formed substitutes for exceptional outcomes.

\subsection{NSwag Client Generation and the Migration Path}
\label{subsec:impl_nswag_migration}

Transitioning to production backends involves leveraging pre-existing generated clients rather than building from scratch. \texttt{program.cs} registers a named \texttt{AntaresGeneratedClient} HTTP client linked to production endpoints via configuration settings, utilising \texttt{AntaresAPIClient} generated by NSwag from OpenAPI specifications, mirroring Anonymisation Service client generation. Basic authentication credentials are pre-configured via user secrets.

Migrating requires implementing \texttt{IMiorelliService} using \texttt{AntaresAPIClient} and updating a single-line dependency injection registration. Tool classes, use cases, and MCP schemas remain unaltered due to abstraction benefits. However, abstract interface wiring represents only initial preparation, as production deployment necessitates writing \texttt{AntaresMiorelliClient}, propagating real caller identities, configuring authentication and role-based access control, and executing integration tests. Antares AI remains a prototype validated against mock backends rather than production environments.
